\documentclass{article}

\newcommand{\doctitle}{Assignment 5} % change document title
\usepackage{ellis_header}

\begin{document}

\section{Chapter 2.3}
\setcounter{problem}{12}
\begin{problem} % Problem 13
Give a combinatorial proof: If $n \ge 1$ and $k \ge 1$, then
\begin{equation*}
    S(n,k) = \sum_{j=1}^{n} \binom{n-1}{j-1} S(n-j, k-1)
\end{equation*}
\end{problem}

\textbf{Answer:}
How many partitions can you make with an $n$-set and $k$-blocks?

\medskip
Prior to starting this proof, I'll provide an example with a $5$-set and $3$ blocks. As a result, the written out formula appears as:

\begin{equation*}
    S(5, 3) = \binom{4}{0} \cdot S(4, 2) + \binom{4}{1} \cdot S(3, 2) + \binom{4}{2} \cdot S(2, 2) + \binom{4}{3} \cdot S(1, 2) + \binom{4}{4} \cdot S(0, 2)
\end{equation*}

For partitions, if $n <k$ then $S(n,k) = 0$. As such, the summation simplfies into.

\begin{equation*}
    S(5, 3) = 1 \cdot S(4, 2) + 4 \cdot S(3, 2) + 6 \cdot S(2, 2)
\end{equation*}

This is equivalent to taking one element, the number $1$ in this example, and anchoring it to be inside a specific partition. At each iteration of $j$ you can add another element to be inside the partition with the anchor value. See the table below for some examples.

\begin{center}
\begin{tabular}{ |c|c|c|c| } 
\hline
j & 1-Block & 2-Block & 3-Block \\
\hline
1 & $\{1\}$ & $\{2\}$ & $\{3, 4, 5\}$ \\
1 & $\{1\}$ & $\{3\}$ & $\{2, 4, 5\}$ \\
1 & $\{1\}$ & $\{4\}$ & $\{2, 3, 5\}$ \\
1 & $\{1\}$ & $\{5\}$ & $\{2, 3, 4\}$ \\ 
1 & $\{1\}$ & $\{2, 3\}$ & $\{4, 5\}$ \\ 
1 & $\{1\}$ & $\{2, 4\}$ & $\{3, 5\}$ \\
1 & $\{1\}$ & $\{2, 5\}$ & $\{3, 4\}$ \\
2 & $\{1,2\}$ & $\{3\}$ & $\{4, 5\}$ \\
2 & $\{1,2\}$ & $\{4\}$ & $\{3, 5\}$ \\
2 & $\{1,2\}$ & $\{5\}$ & $\{3, 4\}$ \\
2 & $\{1,3\}$ & $\{2\}$ & $\{4, 5\}$ \\
2 & $\{1,3\}$ & $\{4\}$ & $\{2, 5\}$ \\
2 & $\{1,3\}$ & $\{5\}$ & $\{2, 4\}$ \\
2 & $\{1,4\}$ & $\{2\}$ & $\{3, 5\}$ \\
2 & $\{1,4\}$ & $\{3\}$ & $\{2, 5\}$ \\
2 & $\{1,4\}$ & $\{5\}$ & $\{2, 4\}$ \\
2 & $\{1,5\}$ & $\{2\}$ & $\{3, 4\}$ \\
2 & $\{1,5\}$ & $\{3\}$ & $\{2, 4\}$ \\
2 & $\{1,5\}$ & $\{4\}$ & $\{2, 3\}$ \\
3 & $\{1,2,3\}$ & $\{4\}$ & $\{5\}$ \\
3 & $\{1,2,4\}$ & $\{3\}$ & $\{5\}$ \\
3 & $\{1,2,5\}$ & $\{3\}$ & $\{4\}$ \\
3 & $\{1,3,4\}$ & $\{2\}$ & $\{5\}$ \\
3 & $\{1,3,5\}$ & $\{2\}$ & $\{4\}$ \\ 
3 & $\{1,4,5\}$ & $\{2\}$ & $\{3\}$ \\  
 \hline
\end{tabular}
\end{center}

LHS: The left hand side counts the possible number of partitions of a n-set, denoted as $[n]$ with k-blocks. This set of possible partitions is $\mathcal{P}$ with size $|\mathcal{P}|$.

\medskip
RHS: Take element $1$ from $[n]$ and anchor it to the first block from $k$. Then $j$ is the number of elements that are inside the first block. This means there are $j-1$ other elments inside the first block and there are $n-1$ options to choose from. This is represented as $\binom{n-1}{j-1}$.

\medskip
Since the first block is defined, there are $k-1$ other blocks to fill in. A selection of $n-j$ elements are available to place in the remaining blocks. This leaves $S(n-j, k-1)$ different sets of partitions.

\medskip
The first block has a unique anchored element inside it which makes it independent of the other $k-1$ partitions. This allows the product rule to be applied for the number of options. $\binom{n-1}{j-1}S(n-j, k-1)$.

\medskip
As $j$ increases in value, the number of elements in the first block grows. As such, each set of partitions with different values of $j$ is disjoint, allowing them to be added without double counting. Represented as $\sum_{j=1}^{n} \binom{n-1}{j-1}S(n-j, k-1)$. This proves that:

\begin{equation*}
    S(n,k) = \sum_{j=1}^{n} \binom{n-1}{j-1} S(n-j, k-1)
\end{equation*}

\section{Chapter 2.4}
\setcounter{problem}{0}
\begin{problem} % Problem 1
You have 40 pieces of candy to distribute among 10 children. Find the number of ways to do this in each of the following situations. Leave your answers in standard notation.
\begin{enumerate}[(a)]
    \item The pieces of candy are different and each child gets at least one piece.
    \item The pieces of candy are the same and each child get any number of pieces.
    \item The pieces of candy are different but you distribute them among 10 of the same paper bags.
    \item The candy is the same you distribute them in 10 bags that are the same and each bag contains at least one piece.
    \item The candy is different and each child gets exactly one pieces, so there is some left over candy.
    \item The candy is different and Frank receives four pieces.
\end{enumerate}
\end{problem}

\textbf{Answer:}
I belive this is as simple as using the table on page 81 of the book for "lookup".

\begin{enumerate}[(a)]
    \item Distribute 40 distinct objects to 10 distinct recipients, each gets at least one. $$S(k,n) \cdot n!$$
    \item Distribute 40 identical objects to 10 distinct reipients, no restrictions. $$\multiset{n}{k}$$
    \item Distribute 40 distinct objects to 10 identical recipients, no restrictions. $$\sum_{i=1}^{n}S(k,i)$$
    \item Distribute 40 identical objects to 10 identical recipients, each gets at least one. $$P(k,n)$$
    \item Distribute 40 distinct objects to 10 distinct recipients, each gets one piece. $$(n)_k$$
    \item Distribute 40 distinct objects to 10 distinct reipients, but Frank gets 4 pieces.$$\binom{40}{4} \cdot 9^{36}$$
\end{enumerate}


\setcounter{problem}{3}
\begin{problem} % Problem 4
Use type vectors to establish the bijection in Theorem 2.4.2.
\begin{thm}
    If $n \ge 1$ and $k \ge 1$, then:
    \begin{equation*}
        P(n,k) = \sum_{j=1}^{k} P(n-k,j)
    \end{equation*}
\end{thm}
\end{problem}

\textbf{Answer:}
Each value inside the original integer partion will have its value decreased by one. For example, take the integer partition of 6 into 3 pieces. This becomes the integer partition of 3 into at most 3 pieces. See the table below for an example.

\begin{center}
\begin{tabular}{ |c|c|c|c| } 
\hline
j & $P(6,3)$ & Subtract 1 & $P(3,j)$ \\
\hline
1 & $\{4,1,1\}$ & $\{4-1,1-1,1-1\}$ & $\{3\}$ \\
2 & $\{3,2,1\}$ & $\{3-1,2-1,1-1\}$ & $\{2, 1\}$ \\
3 & $\{2,2,2\}$ & $\{2-1,2-1,2-1\}$ & $\{1, 1, 1\}$ \\  
 \hline
\end{tabular}
\end{center}

The greatest value possible in the type vector is $n-(k-1)$. This occurs when all the integer partition values are $1$ with the exception of the first value. The example of this in the table being $\{4,1,1\}$. Mathematically represented as $n-(k-1) = n-k+1$. The function takes in the original type vector partition, subtracts 1 from each value, while maintaining the same number of exponents.

\begin{equation*}
f([1^{P_1}2^{P_2}\cdots (n-k+1)^{P_{n-k+1}}]) = [1^{P_2}2^{P_3}\cdots (n-k)^{P_{n-k+1}}]
\end{equation*}

On the left hand side, the integer $n$ and the number of partitions $k$ are:

\begin{equation*}
    n = \sum_{j=1}^{n-k+1} j \cdot P_j \quad \text{and} \quad k = \sum_{j=1}^{n-k+1} P_j
\end{equation*}

On the right hand side, the integer $n-k$ and the number of partitions $k$ are:

\begin{equation*}
    n-k = \sum_{j=2}^{n-k+1} (j-1) \cdot P_j \quad \text{and} \quad k = \sum_{j=2}^{n-k+1} P_j
\end{equation*}

The inverse function from the smaller type vector back to the larger type vector is:

\begin{equation*}
    g([1^{P_1}2^{P_2}\cdots n^{P_{n}}]) = g([1^{n-k}2^{P_1}3^{P_2}\cdots n+1^{P_{n}}])
\end{equation*}

Where:
\begin{equation*}
    k = \sum_{j=1}^n P_j
\end{equation*}


\setcounter{problem}{6}
\begin{problem} % Problem 7
Under what conditions on $n$ and $k$ is the statement $P(n,k) = P(n-1, k-1)$ true?
\end{problem}

\textbf{Answer:} Under the following two options:
\begin{enumerate}
    \item $n=k$ if $n,k \ge 1$
    \item $n=k+1$ if $k > 1$
\end{enumerate}
See the partition number triangle on page 79.


\begin{problem} % Problem 8
Give a bijective proof of the following: The number of partitions of $n$ is equal to the number of partitions of $2n$ into $n$ parts. $P(n) = P(2n, n)$.
\end{problem}

\textbf{Answer:}
The book gives the following hint: \textit{Given a partition of n, add 1 to each existing part and then append enough parts of size 1 to bring the total number of parts to n}. An example will be done for $P(4)$ which translates into $P(8, 4)$.

\begin{center}
\begin{tabular}{ |c|c|c| } 
\hline
$P(4)$ & Add 1 & $P(8,4)$ \\
\hline
$\{4\}$ & $\{4+1, 1, 1, 1\}$ & $\{5,1,1,1\}$ \\
$\{3,1\}$ & $\{3+1,1+1,1, 1\}$ & $\{4, 2, 1, 1\}$ \\
$\{2,2\}$ & $\{2+1,2+1,1,1\}$ & $\{3, 3, 1, 1\}$ \\
$\{1,1,1,1\}$ & $\{1+1,1+1,1+1,1+1\}$ & $\{2,2,2,2\}$ \\
\hline
\end{tabular}
\end{center}

Define $A$ as the set of all partitions of $n$ and $B$ as the set of partitions of $2n$ into $n$ parts. Define the function $f: A \rightarrow B$ by

\begin{equation*}
    f([1^{P_1}2^{P_2}\cdots n^{P_n}]) = [1^{n-k}2^{P_1}3^{P_2}\cdots (n+1)^{P_n}]
\end{equation*}

Where $k = \sum_{j=1}^{n} P_j$. That is, $f$ takes a partition of $n$ adds the value $1$ to all existing parts and adds individual $1$'s to increse the total to $2n$. The inverse of the function $f$ is $g$. Where for each $g(b) = a$ for $f(a) = b$. The function of $g$ is shown below.

\begin{equation*}
    g([1^{P_1}2^{P_2}\cdots (n-k+1)^{P_{n-k+1}}]) = [1^{P_2} 2^{P_3} \cdots (n-k)^{P_{n-k+1}}]
\end{equation*}

\medskip
\textbf{One-to-one:} Let $[1^{p_1}2^{p_2}\cdots n^{p_n}]$ and $[1^{q_1}2^{q_2}\cdots n^{q_n}]$ be partitions in $A$ and assume that 

\begin{equation*}
   f([1^{p_1}2^{p_2}\cdots n^{p_n}]) = f([1^{q_1}2^{q_2}\cdots n^{q_n}]) 
\end{equation*}

This means that $[1^{n-k}2^{p_1}\cdots (n+1)^{p_n}] = [1^{n-k}2^{q_1}\cdots (n+1)^{q_n}]$ and that $p_i = q_i$ for all $i$. Therefore $[1^{p_1}2^{p_2}\cdots n^{p_n}] = [1^{q_1}2^{q_2}\cdots n^{q_n}]$.

\medskip
\textbf{Onto:} Let $[1^{q_1}2^{q_2}3^{q_3} \cdots (n-k+1)^{q_{n-k+1}}]$ be in $B$. Notice that for a parition of $n$ into $k$ integers, the max value of the $k$ integer is $n-k+1$, so we are justified in stopping the type vector there. Moreover

\begin{equation*}
    f([1^{q_2}2^{q_3}\cdots (n-k)^{P_{n-k+1}}]) = [1^{n-k}2^{q_2}3^{q_3} \cdots (n-k+1)^{P_{n-k+1}}]
\end{equation*}

Where $q_1 = n-k$, so $f$ is onto.

\end{document}